\begin{table}[tbp]\centering\small
\caption{Agreement on Announcements and the Contrast in Composition}\label{tab:agree}
\resizebox{\ifdim\width>\textwidth\textwidth\else\width\fi}{!}{%
\begin{tabular}{lrrrrr}
\toprule
FRB/US horizon & Joint response to policy news [90\%] & Normalized equity weight & Normalized dollar weight & Dollar over equity, FRB/US & Ratio of ratios [90\%] \\
\midrule
2 quarters & 0.901 [0.81, 1.01] & $-$0.032 & 0.089 & 2.18 & 2.19 [0.90, 6.01] \\
4 quarters & 0.954 [0.85, 1.07] & $-$0.032 & 0.103 & 2.51 & 2.53 [1.04, 6.91] \\
8 quarters & 0.917 [0.82, 1.03] & $-$0.032 & 0.106 & 2.57 & 2.59 [1.06, 7.08] \\
\bottomrule
\end{tabular}}
\begin{flushleft}\footnotesize\textit{Notes:} FRB/US weights for GDP at the horizon shown, persistent-shift experiment, in funds-rate points; the joint response is the FRB/US weights applied to the responses to a unit of policy news (market index: one). Normalized weights: FRB/US weights divided by that response, so that both indexes read one on a conventional surprise; market weights are 0.609 (10-year yield), $-0.051$ (equity valuation) and 0.065 (dollar). Ratio of ratios: the FRB/US dollar-over-equity ratio divided by the market's (\bl{ag:mkt:ratio}); the interval comes from the bootstrap of the market weights, and the FRB/US weights are held fixed.\end{flushleft}
\end{table}
